\subsection{A Frobenius quadratic bank in characteristic comparable to the square root of length}
\label{sec:scaled-frobenius-padding}

\begin{proposition}\label{prop:scaled-frobenius-padding}
For every prime $p\ge41$, there is a length-$N=9p^2$, dimension-three
Reed--Solomon code over $E=\mathbb F_{p^4}$ and two received words whose
individual and common agreement are at most $2p+3$, with the following
property. At threshold $T=4p$, their affine linear combination has exactly
\[
 M=(p+1)(p-1)^2
\]
exceptional parameters with singleton lists, and one further exceptional
parameter with a list of size two. All other parameters have empty lists.
The threshold is above the DKT first-order curve and below the exact
Johnson threshold. In particular $M=\Theta(N^{3/2})$, the source and common
agreement gaps are at least $2p-3=\Theta(\sqrt N)$, and
$p=\sqrt N/3$ exceeds the message degree.
\end{proposition}

\begin{proof}
Put $B=\mathbb F_{p^2}$ and choose $s\in E^*$ such that
$\eta=s^2\notin B$. Let
\[
 D_0=\{x\ne0:x^2\in B^*\},\qquad D_1=sD_0.
\]
These are disjoint sets of size $2(p^2-1)$. For each $a\in B$ with
$a^{p+1}=1$, let
\[
 I_a=\{y^p-ay:y\in B\}.
\]
The $I_a$ are precisely the $p+1$ one-dimensional $\mathbb F_p$ subspaces
of $B$. On $D_0$ set $f(x)=x^{2p}$ and $g(x)=0$; on $D_1$ set
$f(x)=\eta(x^2/\eta)^p$ and $g(x)=1$.

We use the elementary quadratic classification of the Frobenius word.
On $D_0$, its canonical quadratics are
$Q_{a,b}=aX^2+b$, with $a^{p+1}=1$ and $b\in I_a$.
They have $2p$ matches if $b\ne0$, and $2p-2$ if $b=0$.
Every noncanonical quadratic, including quadratics with coefficients in
$E$, has at most $p$ matches for $p\ge5$.
For completeness, write an arbitrary quadratic as $aX^2+cX+b$.
The core is the union of two $B$-lines $uB^*$, where $u=1$ or
$u=v$ with $v^2$ a nonsquare in $B$. On either line, put $x=ut$ and
\[
 (A,C,B_0)=(a u^{2-2p},c u^{1-2p},b u^{-2p}).
\]
The matching equation is $t^{2p}=At^2+Ct+B_0$.
If these three coefficients are not all in $B$, a $B$-linear functional
on $E$ vanishing on $B$ gives a nonzero polynomial of degree at most two
vanishing at every match on this component. If neither component has
all three coefficients in $B$, there are therefore at most four matches.
If one component does, then $a,b\in B$ and $cu\in B$.
When $c\ne0$, the other component has no matches, because its nonzero
linear term lies outside $B$, whereas all its other terms lie in $B$.
On the coefficientwise $B$-valued component, Frobenius and substitution give
\[
 C^p t^p=H(t),\qquad
 H(t)=(1-A^{p+1})t^2-A^pCt-A^pB_0-B_0^p.
\]
Every match is consequently a zero of
\[
 H(t)^2-C^{2p}(At^2+Ct+B_0).
\]
If this polynomial is nonzero, it has at most four roots. If it is zero,
$H$ has degree at most one, and the nonzero polynomial
$C^pt^p-H(t)$ has at most $p$ roots. Finally, if $c=0$ and a component
is coefficientwise $B$-valued, then $a,b\in B$ on both components.
The substitution $y=x^2$ reduces the entire core to $y^p-ay=b$ on $B^*$,
with two preimages per nonzero $y$. A norm-one $a$ has kernel size $p$
and image $I_a$; a non-norm-one $a$ gives an invertible map. This proves
both the canonical counts and the asserted bound $\max(p,4)=p$ for
all other quadratics.

For the word $f+\lambda g$, the same classification on $D_1$, after
$x=sz$ and division by $\eta$, makes $Q_{a,b}$ canonical precisely when
\[
 b-\lambda=\eta v,\qquad v\in I_a.
\]
Thus a quadratic canonical on both blocks has parameter
$\lambda=b-\eta v$ with $b,v\in I_a$. Since $1,\eta$ are linearly
independent over $B$, the planes $I_a+\eta I_a$ intersect pairwise only
at zero. Every nonzero parameter in their union determines a unique
$a,b,v$. If both $b$ and $v$ are nonzero, the corresponding quadratic
has exactly $4p$ matches; there are $M$ such distinct parameters.
If exactly one is zero it has $4p-2$ matches, and at parameter zero
all doubly canonical quadratics have $4p-4$ matches.
A quadratic canonical on just one block has zero matches on the other:
its norm-one coefficient is already fixed, and the other constant term
fails the relevant image condition. A quadratic canonical on neither
block has at most $2p$ matches in total.

Append $5p^2+4$ distinct points outside $D_0\cup D_1$, avoiding all roots
of $X^3-Q_{a,b}$ for the $p(p+1)$ core canonical quadratics. Put
$f(x)=x^3$, $g(x)=0$ at these new points. This is possible since
\[
 p^4\ge12p^2+3p
\]
for $p\ge5$. Every canonical core quadratic gains zero matches, and
any other quadratic gains at most three. Hence the only parameters
with agreement at least $4p$ are precisely the $M$ parameters above,
each with its unique canonical witness.

Outside the union of the planes, the individual agreement is at most
$2p+3$. The common agreement of $(f,g)$ is also at most $2p+3$.
Indeed, if the proposed direction quadratic is zero, the simultaneous
matches lie in the core and padding and number at most $2p+3$; if it
is one, they lie in $D_1$ and number at most $2p$. Any other direction
quadratic agrees with this zero--one direction at at most four points.
Choose two distinct parameters $\alpha,\beta$ outside the union of
planes, and use the two far words $r=f+\alpha g$ and $s_0=f+\beta g$
as sources. Their common agreement is unchanged by this invertible
change of source coordinates. For $z\ne-1$,
\[
 r+zs_0=(1+z)\left(f+\frac{\alpha+\beta z}{1+z}g\right).
\]
This transports all $M$ singleton exceptions bijectively. The omitted
parameter $\beta$ is not exceptional. At $z=-1$ the word is
$(\alpha-\beta)g$: its list at threshold $4p$ consists exactly of the
two constant polynomials $0$ and $\alpha-\beta$, since every other
quadratic agrees on at most four points.

Finally the exact Johnson threshold is $\sqrt{2N}=\sqrt{18}\,p>4p$.
The finite first-order estimate gives
\[
 N a_1(3/N)\le \sqrt{3N/2}+(3N/8)^{1/4}
 <\frac{15}{4}p+\frac32\sqrt p<4p
\]
for $p\ge41$. All stated conclusions follow.
\end{proof}

Compared with Theorem~\ref{thm:projective-quadratic-line}, this retains
an exceptional population of order $N^{3/2}$ and an agreement gap of
order $\sqrt N$, while increasing the characteristic from order
$N^{1/4}$ to order $\sqrt N$. The ambient field here has size
$p^4=N^2/81$, and the exceptional probability is of order $N^{-1/2}$.
The rate and fractional agreement gap still tend to zero. This is an
extension-field result, not a prime-alphabet construction or a tightness
result for the general DKT bound. The threshold comparison above does
not reuse the preceding construction's finite interpolation certificate.
